% =============================================================================
%  s01_intro.tex — Sección 1: Introducción / Planteamiento del problema
%  Sustituye \lipsum con tu contenido real.
% =============================================================================

\section{Introducción}

% ── Frame: Solo texto — mensaje clave, cita o puente entre secciones ─────────
%    USO: para declaraciones fuertes, preguntas retóricas o transiciones.
%    Copia y adapta. No lleva título de frame ni elementos gráficos.
\begin{frame}[plain]
  \begin{tikzpicture}[remember picture, overlay]
    % Banda lateral de color
    \fill[ColorPrincipal]
      (current page.north west) rectangle
      ([xshift=0.35cm] current page.south west);
    % Línea dorada en el borde de la banda
    \fill[ColorAcento]
      ([xshift=0.32cm] current page.north west) rectangle
      ([xshift=0.35cm] current page.south west);
  \end{tikzpicture}
  \hspace{0.8cm}%
  \begin{minipage}{0.88\textwidth}
    \vspace{2.2cm}
    {\color{ColorAcento}\rule{0.5\textwidth}{1.2pt}}\\[0.5cm]
    {\large\bfseries\color{ColorPrincipal}
      ¿Qué es un Humedal?}\\[0.6cm]
    {\normalsize\color{black}
      De acuerdo con la Convención de Ramsar de 1971, los humedales se definen técnicamente 
      como ecosistemas diversos que actúan como una zona de transición entre los sistemas acuáticos y terrestres, 
      abarcando ambientes tales como turberas, pantanos, llanuras de inundación, lagos y ríos \cite{vergara2024humedales}.
    }\\[0.5cm]
    {\color{ColorAcento}\rule{0.3\textwidth}{0.8pt}}
    \vfill
  \end{minipage}
\end{frame}

% ── Frame: Contexto ───────────────────────────────────────────────────────────
% NOTA \pause: revela bullets uno a uno durante la presentación en vivo,
%   pero genera páginas PDF duplicadas. Descomenta solo si presentas en proyector.
\begin{frame}{Contexto y Problemática}
  Para efectos de este estudio, se establece una diferenciación clave basada en el origen y la gestión de estos sistemas en la capital:
  \begin{itemize}
    \item \textbf{Humedales Naturales (Remanentes lacustres):} Representan los vestigios del antiguo sistema de lagos de la Cuenca de México.
    El caso de \textbf{Xochimilco} destaca como un humedal permanente constituido por una red de canales y lagunas\textit{Sitio Ramsar 1363} \cite{paot2006xochimilco}.
    \item \textbf{Humedales Artificiales y Construidos (Ingeniería Ecológica):} Son sistemas diseñados para emular procesos naturales con fines 
    de remediación. Un ejemplo emblemático es el humedal del \textbf{Bosque de San Juan de Aragón}.
  \end{itemize}
\end{frame}

% ── Frame: Objetivos ─────────────────────────────────────────────────────────
\begin{frame}{Importancia tridimensional en el ámbito urbano}
  \vspace{0.4cm}
  \begin{columns}[T]
    \begin{column}{0.32\textwidth}
      \centering
      {\color{ColorPrincipal}\Large\textbf{Dimensión Social}}\\[0.1cm]
      {\small Ofrecen beneficios directos a la salud pública y el bienestar ciudadano}
      {\small (recreación, investigación y educación ambiental)}
    \end{column}
    \begin{column}{0.32\textwidth}
      \centering
      {\color{ColorPrincipal}\Large\textbf{Dimensión Política}}\\[0.1cm]
      {\small Tensión entre la protección jurídica y la presión del desarrollo urbano}
    \end{column}
    \begin{column}{0.32\textwidth}
      \centering
      {\color{ColorPrincipal}\Large\textbf{Dimensión Ambiental}}\\[0.1cm]
      {\small Los humedales urbanos actúan como reguladores críticos del entorno biofísico}
    \end{column}
  \end{columns}
\end{frame}
